\section{Mirror Constructions}

Fix \(n\ge2\) and put \(N=2n-1\).

\subsection{The standard toric mirror}

Let \(\PP^N\) have homogeneous coordinates \([p_0:\cdots:p_N]\). The standard
toric boundary is
\[
D_{\mathrm{tor}}
=
\sum_{i=0}^N\{p_i=0\}
=
\left\{\prod_{i=0}^N p_i=0\right\},
\]
an anticanonical divisor because it has degree \(N+1\). On the affine chart
\(p_0\ne 0\), write \(x_i=p_i/p_0\) for \(1\leq i\leq N\).

\begin{exam}\label{ex:tor-PN}
For the standard toric structure on \(\PP^N\) \cite{Bat93,Givental,Givental98,HV}, the fan rays are
\(e_1,\ldots,e_N\) and \(-e_1-\cdots-e_N\). After the one-parameter
specialization, the toric superpotential is
\[
W_{\mathrm{tor}}^{(N)}(x_1,\ldots,x_N)
=
x_1+\cdots+x_N+\frac{q}{x_1\cdots x_N}.
\]
The compactifying boundary of this mirror is the toric boundary
\(D_{\mathrm{tor}}\).
\end{exam}

Let \(P_{\mathrm{tor}}^{(N)}=\operatorname{Newt}(W_{\mathrm{tor}}^{(N)})\). Then
\[
P_{\mathrm{tor}}^{(N)}
=
\operatorname{Conv}\{e_1,\ldots,e_N,\,-e_1-\cdots-e_N\}
\subset \RR^N.
\]
With the convention
\[
P^\vee=\{u:\langle u,v\rangle\ge -1\ \text{for all }v\in P\},
\]
its polar dual is
\[
\bigl(P_{\mathrm{tor}}^{(N)}\bigr)^\vee
=
\{u=(u_1,\ldots,u_N)\in\RR^N:
u_i\ge -1\ \text{for all }i,\ 
u_1+\cdots+u_N\le1\}.
\]
